\section{Exact implementation and a conservative integration contract}
\label{sec:implementation}

\subsection{Two producers, separate replay}

The package includes a standard-library Python implementation of both
the ordered-event and distinguished-block algorithms. The first exposes
\texttt{solve}; the second exposes \texttt{solve\_output}. Their independent
checkers expose \texttt{verify\_result} and \texttt{verify\_output}.
They consume the same strict sparse schema:
\begin{lstlisting}[language=Python]
problem = {
    "vertices": 3,
    "columns": 1,
    "edges": [[0, 1], [1, 2], [2, 0]],
    "rows": [[[0, 1]], [[0, 1]], [[0, 1]]],
}
\end{lstlisting}
Row labels are list positions. Each edge is either its two endpoint labels
or \texttt{None}. Counts are positive integers, columns within one row are
distinct, and column labels lie in the declared namespace. Booleans are
rejected where integer labels or counts are required. Counts are bound to
the result by a canonical SHA-256 digest using their original hexadecimal
integer encodings. This avoids mistaking two inputs with identical saturated
matrices but different raw parity for the same authenticated source.

The producer compiles free closure, ordinary reachability, and the
nonlinear supports. Its graphic/partition intersection uses shortest
augmenting paths on cloned edge memberships. This classical primitive is
adapted from the inspected dual-forest-normalization package and is credited
in the source. The new reductions supply overlapping demand groups and
their independent certificate coverage.

Positive replay checks the full spanning tree and literal column leaves
against the supplied matrix. Negative replay independently reconstructs
the logical branch data, uses elementary graphic rank, and validates
complete enumeration coverage. The checker modules import no producer,
SCC optimizer, or matroid-intersection optimizer. Tests deliberately disable
the producer primitives during replay, as well as corrupting tree entries,
pivots, ranks, output assignments, and partition coverage.

\subsection{Result meanings and finite budgets}

\begin{center}
\begin{tabularx}{\linewidth}{@{}lX@{}}
\toprule
Result & Exact meaning \\
\midrule
\texttt{SUCCESS} & A verified spanning tree and full immutable peeling
trace for this input. \\
\texttt{NO\_FOREST} & This input has no successful top tree under the
fixed raw peeling criterion. The attached complete obstruction is checked. \\
\texttt{INCONCLUSIVE} & A stated resource allowance or recursion limit
interrupted discovery; no complete mathematical result is asserted. \\
\bottomrule
\end{tabularx}
\end{center}

The implementation can prove some negatives before exceptional search.
A disconnected eligible graph has no spanning tree. If activation even
from \emph{all} eligible edges leaves a nonempty residue, monotonicity
rules out every tree; that immutable stopping residue is a short obstruction.
These tests are exact and retain their own replay rules.

Structural-work, branch, and partition counters make research caps
explicit. They are not operating-system wall-time or memory limits:
input parsing, hashing, allocations, and big-integer operations have their
ordinary implementation costs. An interrupted search returns
\texttt{INCONCLUSIVE}, not a partial negative certificate. The theoretical
algorithms can stream their searches; the supplied producers retain explicit
negative branches, so actual memory can scale with the fixed-parameter
certificate size.

\subsection{Where to place this operator in the recognizer}

The appropriate insertion point follows authenticated normal-cut geometry
and a chosen lower forest. Preserve the existing inexpensive top choices,
including maximum raw-occurrence weight, because those choices often succeed
without exceptional search. If their immutable peeling tests fail, the
new operator can be used as an exact bounded fallback on the reconstructed
full matrix. The polynomial two-row case remains available. Measured
results in Section~\ref{sec:experiments} determine practical dispatch
priorities; the asymptotic theorem does not make every Python branch faster.

\subsection{A polynomial incumbent prefix retains the FPT guarantee}
\label{sec:bounded-prefix}

The optional \texttt{integration/bounded\_prefix.py} supplies one concrete
dispatcher: run the existing stopping-set search for at most $|R|+1$
visited search nodes, emit a complete result for independent replay if it finishes,
and otherwise invoke \texttt{solve\_output}. The companion dispatcher
checker binds the input digest and calls the appropriate independent
certificate verifier. It does not trust an earlier producer merely
because the earlier implementation is the incumbent.

Each incumbent node performs a polynomial amount of graph and peeling
work, so this prefix has a polynomial total bound. If it completes, its
answer is exact. Otherwise the remaining cost is at most the bound in
Theorem~\ref{thm:factorial-free}. The resulting abstract complete dispatcher therefore
retains $3^kB_{2k+1}\operatorname{poly}(N)$ complexity while allowing easy
cases to finish in the incumbent. A requested cap on the fallback may
still produce \texttt{INCONCLUSIVE}; the full theoretical decision bound
refers to the abstract unrestricted fallback. The Python reference can
still reach its recursion limit and return \texttt{INCONCLUSIVE} even when
no explicit work allowance is supplied.

The particular linear node allowance is a reproducible conservative
choice, not a proved optimum. An obviously oversized disconnected vertex
namespace routes directly to the new sparse certificate producer. This dispatcher is delivered as an integration proposal and
is tested separately from the two new standalone algorithms.

The integration adapter must preserve five pieces of evidence:
\begin{enumerate}
\item the triangulation and normal-coordinate data from which the cut was
reconstructed, including the selected lower forest;
\item the cap-edge map and exterior root, with loops and unavailable rows
represented distinctly;
\item uncapped geometric occurrence counts, including every finite boundary
relation used by the requested completion criterion;
\item the map from dense row labels back to native handles and annular
attachments, and the source-relative charts that legalize cancellations;
\item the original marked boundary information needed by subsequent
hierarchy steps.
\end{enumerate}

After a verified positive matrix result, translate the tree to a rooted
$2/3$ schedule and the peeling trace to $1/2$ cancellations, then replay
them through the existing geometric certificate checker. A finite-matrix
negative should only disable that particular source-relative route.
It must not become a negative unknot verdict. If a cap is reached, continue
the enclosing certified search according to its existing semantics.

The archive deliberately supplies a research module and a precise adapter
contract. It does not silently replace maintained recognizer defaults,
assert that an annular relation is an actual top cap, or claim to have
implemented the missing general marked-homeomorphism compiler.

\subsection{Verification boundary and code limitations}

The independent verifier still shares the input parser and canonical
digest implementation, which form a small explicit trusted base. It checks
the finite mathematical predicate and assumes the caller supplied the
authenticated graph--matrix record. It is not a secure parser for arbitrary
untrusted files under adversarial memory limits, and no such claim is
needed for the present research interface.

The exact weighted theorem and the certified matroid-closure fast path are
proved mathematical extensions. The deliverable implements the two
unweighted forest searches, their complete certificate checks, source
reconstruction studies, and reproducible finite experiments. It contains
no Lean formalization, new theorem about global normal-surface discovery,
or complete quasi-polynomial unknot recognizer.
